\section{Beyond Surface Hopping}\label{sec:beyond_surface_hopping}
While the independent-trajectory methods described in this chapter provide a computationally efficient framework for simulating nonadiabatic dynamics, they are fundamentally limited by the classical treatment of the nuclear degrees of freedom. In particular, these methods have trouble capturing quantum effects such as tunneling, zero-point energy, and wavepacket interference, which can sometimes play a crucial role in the dynamics of molecular systems.

Furthermore, while various decoherence corrections can improve the internal consistency of the \gls{tsh} methods, one might also consider more rigorous approaches that treat the nuclear degrees of freedom quantum mechanically. One notable step beyond the standard \gls{tsh} is the \gls{aims} method.\supercite{10.1021/jp994174i} Unlike the fully variational methods discussed in Sec.~\ref{sec:variational_principles}, \gls{aims} utilizes a basis of frozen Gaussian wavepackets, often called trajectory basis functions. It represents the nuclear wavefunction as a superposition of frozen Gaussian wavepackets whose centers and momenta evolve according to independent classical trajectories rather than coupled variational equations. To capture nonadiabatic transitions, \gls{aims} dynamically spawns new wavepackets on different electronic states when the system passes through regions of strong nonadiabatic coupling. The complex coefficients of this expanding basis are propagated quantum mechanically, allowing \gls{aims} to capture correct wavepacket splitting and interference effects while remaining more computationally tractable than fully coupled methods.

A similar, fully variational approach is the \gls{dd-vmcg} method, which also represents the nuclear wavefunction as a superposition of Gaussian wavepackets.\supercite{10.1039/C8FD00090E} However, unlike \gls{aims}, the \gls{dd-vmcg} method does not rely on a predefined spawning criterion and instead allows the wavepackets to evolve freely according to the time-dependent variational principle. This can lead to a more flexible and accurate description of nonadiabatic dynamics, albeit at a higher computational cost.

Finally, when the highest level of accuracy is required, one might resort to the more advanced \gls{mctdh} method, which represents the nuclear wavefunction as a superposition of time-dependent Hartree products of single-particle functions.\supercite{10.1063/1.472616} While standard grid-based exact quantum dynamics is severely restricted by the exponential scaling of grid points, \gls{mctdh} was specifically designed to alleviate this curse of dimensionality. Although it remains computationally demanding, modern extensions of the theory, such as Multi-Layer \gls{mctdh}, can converge exact or near-exact quantum dynamics for complex systems spanning tens to hundreds of nuclear degrees of freedom.
